\documentclass[a4paper]{scrreprt}

%% Language and font encodings
\usepackage[german]{babel}
\setcounter{secnumdepth}{3} 
\setcounter{tocdepth}{3} 
\usepackage[utf8x]{inputenc}
\usepackage[T1]{fontenc}
\usepackage{courier}

%% Sets page size and margins
\usepackage[a4paper,top=3cm,bottom=2cm,left=3cm,right=3cm,marginparwidth=1.75cm]{geometry}

%% Useful packages
\usepackage{amsmath}
\usepackage{graphicx}
\usepackage[colorinlistoftodos]{todonotes}
\usepackage[colorlinks=true, allcolors=blue, breaklinks = true]{hyperref}
\usepackage{tocstyle}
\usepackage{longtable}
\usetocstyle{standard}
\settocfeature{raggedhook}{\raggedright}
\usepackage{graphicx}
\usepackage{float}
\graphicspath{ {images/} }

\begin{document}	
	\begin{flushright}
		\includegraphics[scale = 0.2]{kit-logo.png}\\[0.5cm]
	\end{flushright}
	\vspace*{2cm}
	
	\begin{center}
		\large Praxis der Softwareentwicklung
		\vspace*{1.5cm}
		
		\textbf{\huge Fridget}
		\vspace*{1cm}
		
		\textbf{\Large Implementierung}
		\vspace*{2cm}
		
		Yunjia Chen, Jasmin Jat, Min Hye Park, Alina Shah, Lisa Wang
		\vspace*{1cm}
		
		\today
		\vspace*{2.5cm}
		
		Betreuung: Erik Burger, Sandro Koch\\[0.5cm]
		IPD\\[0.5cm]
		
		Karlsruher Institut für Technologie
		
	\end{center}
	
	\thispagestyle{empty}
	
	\tableofcontents
	
	\chapter{Einleitung}
	Dies ist das Dokument für die Implementierung der Applikation Fridget – entstanden im Rahmen
	des Softwarepraktikums PSE im Sommersemester 2018.
	Zunächst verlieren wir in diesen Dokument Worte über die Umsetzung von den Muss- und Wunschkriterien unserer App.
	In Abschnitt 3 beschreiben wir die Änderungen gegenüber des Feinentwurfs.
	Die Änderungen gegenüber der Datenstruktur wird in Abschnitt 4 erläutert.
	Ebenso wichtig sind die Unit Tests der Applikation, welche im nachfolgenden Abschnitt 5 aufgelistet sind.
	Um Fachbegriffe zu erläutern, befindet sich im letzten Abschnitt ein Glossar.
	
	\chapter{Umsetzung von Muss- und Wunschkriterien}
	
	
	\chapter{Änderungen gegenüber des Feinentwurfs}
	
	\chapter{Änderungen gegenüber der Datenstruktur}
	
	\chapter{Unit Tests}
	%\documentclass[a4paper]{scrreprt}

%\usepackage[german]{babel}
%\usepackage[utf8]{inputenc}
%\usepackage[T1]{fontenc}
%\usepackage{ae}
%\usepackage{tocbasic}

 \subsection{Package edu.kit.pse.fridget.server.models}
 \subsubsection{\texttt{Class AccessCodeTest}}
 \textbf{Beschreibung} \\
 \textit{Unittest für AccessCode}
 \paragraph*{Methoden}
 \begin{itemize}
    	\item{\texttt{public void buildNew()}}
    	
    	\textit{Testet, ob ein AccessCode-Objekt richtig gebaut wird}
 \end{itemize}

 \subsubsection{\texttt{Class CoolNoteTest}}
 \textbf{Beschreibung} \\
 \textit{Unittest für für CoolNote}
 \paragraph*{Methoden}
 \begin{itemize}
    	\item{\texttt{public void buildForCreation()}}
    	
    	\textit{Testet, ob ein CoolNote-Objekt richtig gebaut wird}
    	
    	\item{\texttt{public void buildForFetching()}}
    	
    	\textit{Testet, ob ein CoolNote-Objekt mit taggedMembershipIds richtig gebaut wird}
 \end{itemize}

 \subsubsection{\texttt{Class DeviceTest}}
 \textbf{Beschreibung} \\
 \textit{Unittest für für Device}
 \paragraph*{Methoden}
 \begin{itemize}
    	\item{\texttt{public void buildNew()}}
    	
    	\textit{Testet, ob ein Device-Objekt richtig gebaut wird}
 \end{itemize}

 \subsubsection{\texttt{Class FlatshareTest}}
 \textbf{Beschreibung} \\
 \textit{Unittest für für Flatshare}
 \paragraph*{Methoden}
 \begin{itemize}
    	\item{\texttt{public String getId()}}
    	
    	\textit{Getter für WG-ID}
 \end{itemize}

 \subsubsection{\texttt{Class MembershipTest}}
 \textbf{Beschreibung} \\
 \textit{Unittest für für Membership}
 \paragraph*{Methoden}
 \begin{itemize}
    	\item{\texttt{public void buildNew()}}
    	
    	\textit{Testet, ob ein Flatshare-Objekt richtig gebaut wird}
 \end{itemize}
 
 \subsubsection{\texttt{Class ReadConfirmationTest}}
 \textbf{Beschreibung} \\
 \textit{Unittest für für ReadConfirmation}
 \paragraph*{Methoden}
 \begin{itemize}
    	\item{\texttt{public void buildNew()}}
    	
    	\textit{Testet, ob ein ReadConfirmation-Objekt richtig gebaut wird}
 \end{itemize}
 
 \subsubsection{\texttt{Class TaggedMemberTest}}
 \textbf{Beschreibung} \\
 \textit{Unittest für TaggedMember}
 \paragraph*{Methoden}
 \begin{itemize}
    	\item{\texttt{public void buildNew()}}
    	
    	\textit{Testet, ob ein TaggedMember-Objekt richtig gebaut wird}
 \end{itemize}
 
 \subsubsection{\texttt{Class UserTest}}
 \textbf{Beschreibung} \\
 \textit{Unittest für User}
 \paragraph*{Methoden}
 \begin{itemize}
    	\item{\texttt{public void buildNew()}}
    	
    	\textit{Testet, ob ein User-Objekt richtig gebaut wird}
 \end{itemize}

 \subsection{Package edu.kit.pse.fridget.server.models.representations}
 \subsubsection{\texttt{Class UserMembershipRepresentationTest}}
 \textbf{Beschreibung} \\
 \textit{Unittest für UserMembershipRepresentation}
 \paragraph*{Methoden}
 \begin{itemize}
    	\item{\texttt{public void buildFromUserAndMembership()}}
    	
    	\textit{Testet, ob ein UserMembershipRepresentation-Objekt richtig gebaut wird}
 \end{itemize}
 
 \newpage
 
 \subsection{Package edu.kit.pse.fridget.server.services}
 \subsubsection{\texttt{Class AccessCodeServiceTest}}
 \textbf{Beschreibung} \\
 \textit{Unittest für AccessCodeService}
 \paragraph*{Methoden}
 \begin{itemize}
    	\item{\texttt{public void generateAccessCode()}}
    	
    	\textit{Testet, ob ein Zugangscode richtig generiert wird}
    	
    	\item{\texttt{public void generateAccessCode\_WithIncorrectFlatshareId()}}
    	
    	\textit{Testet, ob eine EntityUnprocessableException geworfen wird, falls die gegebene flatshareId falsch ist}
 \end{itemize}
 
 \subsubsection{\texttt{Class CoolNoteServiceTest}}
 \textbf{Beschreibung} \\
 \textit{Unittest für CoolNoteService}
\paragraph*{Methoden}
\begin{itemize}
      	\item{\texttt{public void getAllCoolNotes\_WithIncorrectFlatshareId()}}
      	
      	\textit{Testet, ob eine EntityNotFoundException geworfen wird, falls die gegebene flatshareId falsch ist}
      	
      	\item{\texttt{public void getCoolNote\_WithIncorrectId}}
      	
      	\textit{Testet, ob eine EntityNotFoundException geworfen wird, falls die gegebene coolNoteId falsch ist}
      	
      	\item{\texttt{public void saveCoolNote\_WithIncorrectCreatorMembershipId()}}
      	
      	\textit{Testet, ob eine EntityUnprocessableException geworfen wird, falls die gegebene creatorMembershipId falsch ist}
      	
      	\item{\texttt{public void saveCoolNote\_WithIncorrectPosition}}
      	
      	\textit{Testet, ob eine EntityConflictException geworfen wird, falls die gegebene position falsch ist}
      	
      	\item{\texttt{public void deleteCoolNote\_WithIncorrectId}}
      	
      	\textit{Testet, ob eine EntityConflictException geworfen wird, falls die gegebene coolNoteId falsch ist}
\end{itemize}

 \subsubsection{\texttt{Class DeviceServiceTest}}
 \textbf{Beschreibung} \\
 \textit{Unittest für DeviceService}
 \paragraph*{Methoden}
 \begin{itemize}
    	\item{\texttt{public void saveDevice\_WithIncorrectUserId()}}
    	
    	\textit{Testet, ob eine EntityUnprocessableException geworfen wird, falls die gegebene userId falsch ist}

    	\item{\texttt{public void updateDevice\_WithIncorrectUserId()}}
    	
    	\textit{Testet, ob eine EntityUnprocessableException geworfen wird, falls die gegebene userId falsch ist}
 \end{itemize}

 \subsubsection{\texttt{Class FlatshareServiceTest}}
 \textbf{Beschreibung} \\
 \textit{Unittest für Flatshare}
 \paragraph*{Methoden}
 \begin{itemize}
    	\item{\texttt{public void getFlatshare\_WithIncorrectId}}
    	
    	\textit{Testet, ob eine EntityNotFoundException geworfen wird, falls die gegebene flatshareId falsch ist}
    	
    	\item{\texttt{public void saveFlatshare\_WithIncorrectUserId}}
    	
    	\textit{Testet, ob eine EntityUnprocessableException geworfen wird, falls die gegebene userId falsch ist}
 \end{itemize}
 
 \subsubsection{\texttt{Class FrozenNoteServiceTest}}
 \textbf{Beschreibung} \\
 \textit{Unittest für FrozenNote}
 \paragraph*{Methoden}
 \begin{itemize}
    	\item{\texttt{public void getAllFrozenNotes\_WithIncorrectFlatshareId()}}
    	
    	\textit{Testet, ob eine EntityNotFoundException geworfen wird, falls die gegebene flatshareId falsch ist}
    	
    	\item{\texttt{public void getFrozenNote\_WithIncorrectId()}}
    	
    	\textit{Testet, ob eine EntityNotFoundException geworfen wird, falls die gegebene frozenNoteId falsch ist}
    	
    	\item{\texttt{public void updateFrozenNote\_WithIncorrectFlatshareId()}}
    	
    	\textit{Testet, ob eine EntityUnprocessableException geworfen wird, falls die gegebene flatshareId falsch ist}

    	\item{\texttt{public void updateFrozenNote\_WithIncorrectPosition()}}
    	
    	\textit{Testet, ob eine EntityConflictException geworfen wird, falls die gegebene position falsch ist}
 \end{itemize}
 
 \subsubsection{\texttt{Class MembershipServiceTest}}
 \textbf{Beschreibung} \\
 \textit{Unittest für MembershipService}
 \paragraph*{Methoden}
 \begin{itemize}
    	\item{\texttt{public void getAllMembers\_WithIncorrectFlatshareId()}}
    	
    	\textit{Testet, ob eine EntityNotFoundException geworfen wird, falls die gegebene flatshareId falsch ist}
    	
    	\item{\texttt{public void getMember\_WithIncorrectFlatshareId()}}
    	
    	\textit{Testet, ob eine EntityNotFoundException geworfen wird, falls die gegebene flatshareId falsch ist}
    	
    	\item{\texttt{public void saveMembership\_WithIncorrectAccessCode()}}
    	
    	\textit{Testet, ob eine EntityNotFoundException geworfen wird, falls der gegebene accessCode falsch ist}
    	
    	\item{\texttt{public void saveMembership\_WithIncorrectUserId()}}
    	
    	\textit{Testet, ob eine EntityUnprocessableException geworfen wird, falls die gegebene userId falsch ist}
    	
    	\item{\texttt{public void saveMembership\_WithFullFlatshare()}}
    	
    	\textit{Testet, ob eine EntityConflictException geworfen wird, falls die WG schon voll ist}
    	
    	\item{\texttt{public void deleteMembership\_WithIncorrectFlatshareId()}}
    	
    	\textit{Testet, ob eine EntityConflictException geworfen wird, falls die gegebene flatshareId falsch ist}
    	
    	\item{\texttt{public void deleteMembership\_WithIncorrectUserId()}}
    	
    	\textit{Testet, ob eine EntityConflictException geworfen wird, falls die gegebene userId falsch ist}
 \end{itemize}
 
 \subsubsection{\texttt{Class ReadConfirmationServiceTest}}
 \textbf{Beschreibung} \\
 \textit{Unittest für ReadConfirmationService}
 \paragraph*{Methoden}
 \begin{itemize}
    	\item{\texttt{public void getAllMemberships\_WithIncorrectCoolNoteId()}}
    	
    	\textit{Testet, ob eine EntityNotFoundException geworfen wird, falls die gegebene coolNoteId falsch ist}
    	
    	\item{\texttt{public void saveReadConfirmation\_WithIncorrectCoolNoteId()}}
    	
    	\textit{Testet, ob eine EntityUnprocessableException geworfen wird, falls die gegebene coolNoteId falsch ist}
    	
    	\item{\texttt{public void saveReadConfirmation\_WithIncorrectMembershipId()}}
    	
    	\textit{Testet, ob eine EntityUnprocessableException geworfen wird, falls die gegebene membershipId falsch ist}
    	
    	\item{\texttt{public void saveReadConfirmation\_WithExistedCoolNoteIdAndMembershipId()}}
    	
    	\textit{Testet, ob eine EntityConflictException geworfen wird, falls die gegebenen coolNoteId und membershipId schon existieren}
    	
    	\item{\texttt{public void deleteReadConfirmation\_WithIncorrectCoolNoteId()}}
    	
    	\textit{Testet, ob eine EntityConflictException geworfen wird, falls die gegebene coolNoteId falsch ist} 
    	
    	\item{\texttt{public void deleteReadConfirmation\_WithIncorrectMembershipId()}}
    	
    	\textit{Testet, ob eine EntityConflictException geworfen wird, falls die gegebene membershipId falsch ist}
 \end{itemize}

 \subsubsection{\texttt{Class MagnetColorServiceTest}}
 \textbf{Beschreibung} \\
 \textit{Unittest für MagnetColorService}
\paragraph*{Methoden}
\begin{itemize}
	\item{\texttt{public void getRandomColor()}}
	
	\textit{Testet, ob eine zufällige Farbe generiert wird}
	
	\item{\texttt{public void getAvailableRandomColor()}}
	
	\textit{Testet, ob eine verfügbare zufällige Farbe generiert wird}	
\end{itemize}

\subsubsection{\texttt{Class JwtServiceTest}}
\textbf{Beschreibung} \\
\textit{Unittest für JwtService}
\paragraph*{Methoden}
\begin{itemize}
   	\item{\texttt{public void encodeAndDecode\_ReturnsCorrectUserId()}}
   	
   	\textit{Testet, ob ein JWT richtig encodiert und decodiert wird}
\end{itemize}

\newpage

 \subsection{Package edu.kit.pse.fridget.server.controllers}
 \subsubsection{\texttt{Class AccessCodeControllerIntegrationTest}}
 \textbf{Beschreibung} \\
 \textit{Integration Test für AccessCodeController}
 \paragraph*{Methoden}
 \begin{itemize}
 	\item{\texttt{public void generateAccessCode()}}
 	
 	\textit{Testet, ob ein Zugangscode als ResponseEntity mit einem entsprechenden Statuscode zurückgegeben wird}
 	
 	\item{\texttt{public void generateAccessCode\_WithIncorrectFlatshareId\_Returns\-UnprocessableEntity()}}
 	
 	\textit{Testet, ob eine Fehlermeldung mit einem entsprechenden Statuscode im Response zurückgegeben wird, falls die gegebene flatshareId falsch ist}
 \end{itemize}
 
 \subsubsection{\texttt{Class CoolNoteControllerIntegrationTest}}
 \textbf{Beschreibung} \\
 \textit{Integration Test für CoolNoteController}
 \paragraph*{Methoden}
 \begin{itemize}
 	\item{\texttt{public void getAllCoolNotes()}}
 	
 	\textit{Testet, ob alle Cool Notes als ResponseEntity mit einem entsprechenden Statuscode zurückgegeben werden}
 	
 	\item{\texttt{public void getAllCoolNotes\_WithIncorrectFlatshareId()}}
 	
 	\textit{Testet, ob eine Fehlermeldung mit einem entsprechenden Statuscode im Response zurückgegeben wird, falls die gegebene flatshareId falsch ist}
 	
 	\item{\texttt{public void getCoolNote()}}
 	
 	\textit{Testet, ob eine Cool Note als ResponseEntity mit einem entsprechenden Statuscode zurückgegeben wird}
 	
 	\item{\texttt{public void getCoolNote\_WithIncorrectId()}}
 	
 	\textit{Testet, ob eine Fehlermeldung mit einem entsprechenden Statuscode im Response zurückgegeben wird, falls die gegebene coolNoteId falsch ist}
 	
 	\item{\texttt{public void saveCoolNote()}}
 	
 	\textit{Testet, ob eine Cool Note als ResponseEntity mit einem entsprechenden Statuscode zurückgegeben wird}
 	
 	\item{\texttt{public void saveCoolNote\_WithIncorrectCreatorMembershipId\_Returns\-UnprocessableEntity()}}
 	
 	\textit{Testet, ob eine Fehlermeldung mit einem entsprechenden Statuscode im Response zurückgegeben wird, falls die gegebene creatorMembershipId falsch ist}
 	
 	\item{\texttt{public void saveCoolNote\_WithIncorrectPosition\_ReturnsConflict()}}
 	
 	\textit{Testet, ob eine Fehlermeldung mit einem entsprechenden Statuscode im Response zurückgegeben wird, falls die gegebene position falsch ist}
 	
 	\item{\texttt{public void deleteCoolNote()}}
 	
 	\textit{Testet, ob ein entsprechender Statuscode im Response zurückgegeben wird}
 	
 	\item{\texttt{public void deleteCoolNote\_WithIncorrectId\_ReturnsConflict}}
 	
 	\textit{Testet, ob eine Fehlermeldung mit einem entsprechenden Statuscode im Response zurückgegeben wird, falls die gegebene coolNoteId falsch ist}
 \end{itemize}
 
 \subsubsection{\texttt{Class DeviceControllerIntegrationTest}}
 \textbf{Beschreibung} \\
 \textit{Integration Test für DeviceController}
 \paragraph*{Methoden}
 \begin{itemize}
 	\item{\texttt{public void saveDevice()}}
 	
 	\textit{Testet, ob ein Device als ResponseEntity mit einem entsprechenden Statuscode zurückgegeben wird}
 	
 	\item{\texttt{public void saveDevice\_WithIncorrectUserId()}}
 	
 	\textit{Testet, ob eine Fehlermeldung mit einem entsprechenden Statuscode im Response zurückgegeben wird, falls die gegebene userId falsch ist}
 	
 	\item{\texttt{public void updateDevice()}}
 	
 	\textit{Testet, ob ein Device als ResponseEntity mit einem entsprechenden Statuscode zurückgegeben wird}
 	
 	\item{\texttt{public void updateDevice\_WithIncorrectUserId\_Returns\-UnprocessableEntity()}}
 	
 	\textit{Testet, ob eine Fehlermeldung mit einem entsprechenden Statuscode im Response zurückgegeben wird, falls die gegebene userId falsch ist}
 \end{itemize}
 
 \subsubsection{\texttt{Class FlatshareControllerIntegrationTest}}
 \textbf{Beschreibung} \\
 \textit{Integration Test für FlatshareController}
 \paragraph*{Methoden}
 \begin{itemize}
 	\item{\texttt{public void getFlatshare()}}
 	
 	\textit{Testet, ob eine Flatshare als ResponseEntity mit einem entsprechenden Statuscode zurückgegeben wird}
 	
 	\item{\texttt{public void getFlatshre\_WithIncorrectId\_ReturnsNotFound()}}
 	
 	\textit{Testet, ob eine Fehlermeldung mit einem entsprechenden Statuscode im Response zurückgegeben wird, falls die gegebene flatshareId falsch ist}
 	
 	\item{\texttt{public void saveFlatshare()}}
 	
 	\textit{Testet, ob eine Flatshare als ResponseEntity mit einem entsprechenden Statuscode zurückgegeben wird}
 	
 	\item{\texttt{public void saveFlatshare\_WithIncorrectUserId\_Returns\-UnprocessableEntity}}
 	
 	\textit{Testet, ob eine Fehlermeldung mit einem entsprechenden Statuscode im Response zurückgegeben wird, falls die gegebene userId falsch ist}
 \end{itemize}
 
 \subsubsection{\texttt{Class FrozenNoteControllerIntegrationTest}}
 \textbf{Beschreibung} \\
 \textit{Integration Test für FrozenNoteController}
 \paragraph*{Methoden}
 \begin{itemize}
 	\item{\texttt{public void getAllFrozenNotes()}}
 	
 	\textit{Testet, ob alle Frozen Notes als ResponseEntity mit einem entsprechenden Statuscode zurückgegeben werden}
 	
 	\item{\texttt{public void getAllFrozenNotes\_WithIncorrectFlatshareId()}}
 	
 	\textit{Testet, ob eine Fehlermeldung mit einem entsprechenden Statuscode im Response zurückgegeben wird, falls die gegebene flatshareId falsch ist}
 	
 	\item{\texttt{public void getFrozenNote()}}
 	
 	\textit{Testet, ob eine Frozen Note als ResponseEntity mit einem entsprechenden Statuscode zurückgegeben wird}
 	
 	\item{\texttt{public void getFrozenNote\_WithIncorrectId()}}
 	
 	\textit{Testet, ob eine Fehlermeldung mit einem entsprechenden Statuscode im Response zurückgegeben wird, falls die gegebene frozenNoteId falsch ist}
 	
 	\item{\texttt{public void updateFrozenNote()}}
 	
 	\textit{Testet, ob eine Frozen Note als ResponseEntity mit einem entsprechenden Statuscode zurückgegeben wird}
 	
 	\item{\texttt{public void updateFrozenNote\_WithIncorrectFlatshareId\_Returns\-UnprocessableEntity()}}
 	
 	\textit{Testet, ob eine Fehlermeldung mit einem entsprechenden Statuscode im Response zurückgegeben wird, falls die gegebene flatshareId falsch ist}
 	
 	\item{\texttt{public void updateFrozenNote\_WithIncorrectPosition\_ReturnsConflict()}}
 	
 	\textit{Testet, ob eine Fehlermeldung mit einem entsprechenden Statuscode im Response zurückgegeben wird, falls die gegebene position falsch ist}
 \end{itemize}
 
 \subsubsection{\texttt{Class MembershipControllerIntegrationTest}}
 \textbf{Beschreibung} \\
 \textit{Integration Test für MembershipController}
 \paragraph*{Methoden}
 \begin{itemize}
 	\item{\texttt{public void getAllMembers()}}
 	
 	\textit{Testet, ob alle Members als ResponseEntity mit einem entsprechenden Statuscode zurückgegeben werden}
 	
 	\item{\texttt{public void getAllMembers\_WithIncorrectFlatshareId()}}
 	
 	\textit{Testet, ob eine Fehlermeldung mit einem entsprechenden Statuscode im Response zurückgegeben wird, falls die gegebene flatshareId falsch ist}
 	
 	\item{\texttt{public void getMember()}}
 	
 	\textit{Testet, ob ein Member als ResponseEntity mit einem entsprechenden Statuscode zurückgegeben wird}
 	
 	\item{\texttt{public void getMember\_WithIncorrectUserId()}}
 	
 	\textit{Testet, ob eine Fehlermeldung mit einem entsprechenden Statuscode im Response zurückgegeben wird, falls die gegebene userId falsch ist}
 	
 	\item{\texttt{public void getMember\_WithIncorrectFlatshareId()}}
 	
 	\textit{Testet, ob eine Fehlermeldung mit einem entsprechenden Statuscode im Response zurückgegeben wird, falls die gegebene flatshareId falsch ist}
 	
 	\item{\texttt{public void saveMembership\_WithAccessCodeDeleted()}}
 	
 	\textit{Testet, ob ein Member als ResponseEntity mit einem entsprechenden Statuscode zurückgegeben wird und der Zugangscode gelöscht wird}
 	
 	\item{\texttt{public void saveMembership\_WithIncorrectAccessCode\_ReturnsNotFound()}}
 	
 	\textit{Testet, ob eine Fehlermeldung mit einem entsprechenden Statuscode im Response zurückgegeben wird, falls der gegebene Zugangscode falsch ist}
 	
 	\item{\texttt{public void saveMembership\_WithIncorrectUserId\_Returns\-UnprocessableEntity()}}
 	
 	\textit{Testet, ob eine Fehlermeldung mit einem entsprechenden Statuscode im Response zurückgegeben wird, falls die gegebene userId falsch ist}
 	
 	\item{\texttt{public void saveMembership\_WithExistedUserId\_ReturnsConflict()}}
 	
 	\textit{Testet, ob eine Fehlermeldung mit einem entsprechenden Statuscode im Response zurückgegeben wird, falls das Member schon existiert}
 	
 	\item{\texttt{public void deleteMembership()}}
 	
 	\textit{Testet, ob ein entsprechender Statuscode im Response zurückgegeben wird}
 	
 	\item{\texttt{public void deleteMembership\_WithIncorrectFlatshareId\_ReturnsConflict()}}
 	
 	\textit{Testet, ob eine Fehlermeldung mit einem entsprechenden Statuscode im Response zurückgegeben wird, falls die gegebene flatshareId falsch ist}
 	
 	\item{\texttt{public void deleteMembership\_WithIncorrectUserId\_ReturnsConflict()}}
 	
 	\textit{Testet, ob eine Fehlermeldung mit einem entsprechenden Statuscode im Response zurückgegeben wird, falls die gegebene userId falsch ist}
 \end{itemize}
 
 \subsubsection{\texttt{Class ReadConfirmationControllerIntegrationTest}}
 \textbf{Beschreibung} \\
 \textit{Integration Test für ReadConfirmationController}
 \paragraph*{Methoden}
 \begin{itemize}
 	\item{\texttt{public void getAllMemberships()}}
 	
 	\textit{Testet, ob alle Members als ResponseEntity mit einem entsprechenden Statuscode zurückgegeben werden}
 	
 	\item{\texttt{public void getAllMemberships\_WithIncorrectCoolNoteId()}}
 	
 	\textit{Testet, ob eine Fehlermeldung mit einem entsprechenden Statuscode im Response zurückgegeben wird, falls die gegebene coolNoteId falsch ist}
 	
 	\item{\texttt{public void saveReadConfirmation()}}
 	
 	\textit{Testet, ob eine Lesebestätigung als ResponseEntity mit einem entsprechenden Statuscode zurückgegeben wird}
 	
 	\item{\texttt{public void saveReadConfirmation\_WithIncorrectMembershipId\_Returns\-UnprocessableEntity()}}
 	
 	\textit{Testet, ob eine Fehlermeldung mit einem entsprechenden Statuscode im Response zurückgegeben wird, falls die gegebene membershipId falsch ist}
 	
 	\item{\texttt{public void saveReadConfirmation\_WithIncorrectCoolNoteId\_Returns\-UnprocessableEntity()}}
 	
 	\textit{Testet, ob eine Fehlermeldung mit einem entsprechenden Statuscode im Response zurückgegeben wird, falls die gegebene coolNoteId falsch ist}
 	
 	\item{\texttt{public void saveReadConfirmation\_WithExistedCoolNoteId\-AndMembershipId\_ReturnsConflict()}}
 	
 	\textit{Testet, ob eine Fehlermeldung mit einem entsprechenden Statuscode im Response zurückgegeben wird, falls die gegebenen coolNoteId und membershipId schon existieren} 	
 	
 	\item{\texttt{public void deleteReadConfirmation()}}
 	
 	\textit{Testet, ob ein entsprechender Statuscode im Response zurückgegeben wird}
 	
 	\item{\texttt{public void deleteReadConfirmation\_WithIncorrectCoolNoteId\_Returns\-Conflict()}}
 	
 	\textit{Testet, ob eine Fehlermeldung mit einem entsprechenden Statuscode im Response zurückgegeben wird, falls die gegebene coolNoteId falsch ist}
 	
 	\item{\texttt{public void deleteReadConfirmation\_WithIncorrectMembershipId\_Returns\-Conflict()}}
 	
 	\textit{Testet, ob eine Fehlermeldung mit einem entsprechenden Statuscode im Response zurückgegeben wird, falls die gegebene membershipId falsch ist} 	
 \end{itemize}
 
 \subsubsection{\texttt{Class UserControllerIntegrationTest}}
 \textbf{Beschreibung} \\
 \textit{Integration Test für UserController}
 \paragraph*{Methoden}
 \begin{itemize}
 	\item{\texttt{public void registerOrLogin\_WithIdToken()}}
 	
 	\textit{Testet, ob ein Benutzer als ResponseEntity mit einem entsprechenden Statuscode zurückgegeben wird}
 \end{itemize}
	
	\chapter{Glossar}
	
	\chapter{Anhang}
\end{document}